%=====================================
%   20260204-rits.tex
%   Microlocal Morse Theory and Truncated Microsupport
%   toshi2019
%=====================================
\documentclass[dvipdfmx,10pt,aspectratio=169,leqno]{beamer}% dvipdfmxしたい
\title{Microlocal Morse Theory and Truncated Microsupport\\
（超局所モース理論と打切り超局所台）
\subtitle{修論審査会}}
\date{2026年2月4日}
\author{大柴 寿浩}

\institute[NITech]
{
北海道大学大学院理学院数学専攻
}


\usetheme{metropolis}
%\usetheme{Singapore}
\usecolortheme{rose}
\usefonttheme{professionalfonts}
\setbeamertemplate{navigaton symbols}{\false}
\setbeamertemplate{footline}[frame number]
%\usepackage{graphicx,xcolor}





\usepackage{bxdpx-beamer}% dvipdfmxなので必要
\usepackage{pxjahyper}% 日本語で'しおり'したい
\usepackage{minijs}% min10ヤダ
\renewcommand{\kanjifamilydefault}{\gtdefault}
\renewcommand{\emph}[1]{{\upshape\bfseries #1}}

\usepackage{amsmath}
\usepackage{amsthm}
\usepackage{tikz}
\usepackage{color}
\usepackage{ascmac}
\usepackage{amsfonts}
\usepackage{mathrsfs}
\usepackage[mathscr]{eucal}
\usepackage{mathtools}
\usepackage{amssymb}
\usepackage{graphicx}
\usepackage{fancybox}
\usepackage{enumerate}
\usepackage{verbatim}
\usepackage{subfigure}
\usepackage{proof}
\usepackage{listings}
\usepackage{otf}
\usepackage[all]{xy}
\usepackage{amscd}
%\usepackage[dvipdfmx]{hyperref}

\usepackage{xcolor}
\definecolor{darkgreen}{rgb}{0,0.45,0} 
\definecolor{darkred}{rgb}{0.75,0,0}
\definecolor{darkblue}{rgb}{0,0,0.6} 
\hypersetup{
    colorlinks=true,
    citecolor=darkgreen,
    linkcolor=darkblue,
    urlcolor=darkblue,
}

\usetikzlibrary{positioning}


\usepackage{latexsym}
\usepackage{wrapfig}
\usepackage{layout}
\usepackage{url}

\usepackage{okumacro}

\usepackage{comment}

%\usepackage{jdash}
%\usepackage{pxjahyper}
% ----------------------------
% commmand
% ----------------------------
% 執筆に便利なコマンド集です. 
% コマンドを追加する場合は下のスペースへ. 

% 集合の記号 (黒板文字)
\newcommand{\NN}{\mathbb{N}}
\newcommand{\ZZ}{\mathbb{Z}}
\newcommand{\QQ}{\mathbb{Q}}
\newcommand{\RR}{\mathbb{R}}
\newcommand{\CC}{\mathbb{C}}
\newcommand{\PP}{\mathbb{P}}
\newcommand{\KK}{\mathbb{K}}


% 集合の記号 (太文字)
\newcommand{\nn}{\mathbf{N}}
\newcommand{\zz}{\mathbf{Z}}
\newcommand{\qq}{\mathbf{Q}}
\newcommand{\rr}{\mathbf{R}}
\newcommand{\cc}{\mathbf{C}}
\newcommand{\pp}{\mathbf{P}}
\newcommand{\kk}{\mathbf{k}}

% 特殊な写像の記号
\newcommand{\ev}{\mathop{\mathrm{ev}}\nolimits} % 値写像
\newcommand{\pr}{\mathop{\mathrm{pr}}\nolimits} % 射影

% スクリプト体にするコマンド
%   例えば　{\mcal C} のように用いる
\newcommand{\mcal}{\mathcal}

% 花文字にするコマンド 
%   例えば　{\h C} のように用いる
\newcommand{\h}{\mathscr}

% ヒルベルト空間などの記号
\newcommand{\F}{\mcal{F}}
\newcommand{\X}{\mcal{X}}
\newcommand{\Y}{\mcal{Y}}
\newcommand{\Hil}{\mcal{H}}
\newcommand{\RKHS}{\Hil_{k}}
\newcommand{\Loss}{\mcal{L}_{D}}
\newcommand{\MLsp}{(\X, \Y, D, \Hil, \Loss)}

% 偏微分作用素の記号
\newcommand{\p}{\partial}

% 角カッコの記号 (内積は下にマクロがあります)
\newcommand{\lan}{\langle}
\newcommand{\ran}{\rangle}



% 圏の記号など
\newcommand{\Set}{{\bf Set}}
\newcommand{\Vect}{{\bf Vect}}
\newcommand{\FDVect}{{\bf FDVect}}
%\newcommand{\Ring}{{\bf Ring}}
\newcommand{\Ab}{{\mathsf{Ab}}}
\newcommand{\Mod}{\mathop{\mathrm{Mod}}\nolimits}
\newcommand{\Modf}{\mathop{\mathrm{Mod}^\mathrm{f}}\nolimits}
\newcommand{\CGA}{{\bf CGA}}
\newcommand{\GVect}{{\bf GVect}}
\newcommand{\Lie}{{\bf Lie}}
\newcommand{\dLie}{{\bf Liec}}



% 射の集合など
\newcommand{\Map}{\mathop{\mathrm{Map}}\nolimits} % 写像の集合
\newcommand{\Hom}{\mathop{\mathrm{Hom}}\nolimits} % 射集合
\newcommand{\End}{\mathop{\mathrm{End}}\nolimits} % 自己準同型の集合
\newcommand{\Aut}{\mathop{\mathrm{Aut}}\nolimits} % 自己同型の集合
\newcommand{\Mor}{\mathop{\mathrm{Mor}}\nolimits} % 射集合
\newcommand{\Ker}{\mathop{\mathrm{Ker}}\nolimits} % 核
\newcommand{\Img}{\mathop{\mathrm{Im}}\nolimits} % 像
\newcommand{\Cok}{\mathop{\mathrm{Coker}}\nolimits} % 余核
\newcommand{\Cim}{\mathop{\mathrm{Coim}}\nolimits} % 余像

% その他便利なコマンド
\newcommand{\dip}{\displaystyle} % 本文中で数式モード
\newcommand{\e}{\varepsilon} % イプシロン
\newcommand{\dl}{\delta} % デルタ
\newcommand{\pphi}{\varphi} % ファイ
\newcommand{\ti}{\tilde} % チルダ
\newcommand{\pal}{\parallel} % 平行
\newcommand{\op}{{\rm op}} % 双対を取る記号
\newcommand{\lcm}{\mathop{\mathrm{lcm}}\nolimits} % 最小公倍数の記号
\newcommand{\Probsp}{(\Omega, \F, \P)} 
\newcommand{\argmax}{\mathop{\rm arg~max}\limits}
\newcommand{\argmin}{\mathop{\rm arg~min}\limits}





% ================================
% コマンドを追加する場合のスペース 
%\newcommand{\OO}{\mcal{O}}



\makeatletter
\renewenvironment{proof}[1][\proofname]{\par
  \pushQED{\qed}%
  \normalfont \topsep6\p@\@plus6\p@\relax
  \trivlist
  \item[\hskip\labelsep
%        \itshape
         \bfseries
%    #1\@addpunct{.}]\ignorespaces
    {#1}]\ignorespaces
}{%
  \popQED\endtrivlist\@endpefalse
}
\makeatother

\renewcommand{\proofname}{\textrm{Proof.}}



%\renewcommand\proofname{\bf 証明} % 証明
\numberwithin{equation}{subsection}
\newcommand{\cTop}{\textsf{Top}}
%\newcommand{\cOpen}{\textsf{Open}}
\newcommand{\Op}{\mathop{\textsf{Op}}\nolimits}
\newcommand{\Ob}{\mathop{\textrm{Ob}}\nolimits}
\newcommand{\id}{\mathop{\mathrm{id}}\nolimits}
\newcommand{\pt}{\mathop{\mathrm{pt}}\nolimits}
\newcommand{\res}{\mathop{\rho}\nolimits}
\newcommand{\A}{\mcal{A}}
\newcommand{\B}{\mcal{B}}
\newcommand{\C}{\mcal{C}}
\newcommand{\D}{\mcal{D}}
\newcommand{\E}{\mcal{E}}
\newcommand{\G}{\mcal{G}}
%\newcommand{\H}{\mcal{H}}
\newcommand{\I}{\mcal{I}}
\newcommand{\J}{\mcal{J}}
\newcommand{\OO}{\mcal{O}}
\newcommand{\Ring}{\mathop{\textsf{Ring}}\nolimits}
\newcommand{\cAb}{\mathop{\textsf{Ab}}\nolimits}
%\newcommand{\Ker}{\mathop{\mathrm{Ker}}\nolimits}
\newcommand{\im}{\mathop{\mathrm{Im}}\nolimits}
\newcommand{\Coker}{\mathop{\mathrm{Coker}}\nolimits}
\newcommand{\Coim}{\mathop{\mathrm{Coim}}\nolimits}
\newcommand{\rank}{\mathop{\mathrm{rank}}\nolimits}
\newcommand{\Ht}{\mathop{\mathrm{Ht}}\nolimits}
\newcommand{\supp}{\mathop{\mathrm{supp}}\nolimits}
\newcommand{\colim}{\mathop{\mathrm{colim}}}
\newcommand{\Tor}{\mathop{\mathrm{Tor}}\nolimits}

\newcommand{\cat}{\mathscr{C}}

%筆記体
\newcommand{\cA}{\mcal{A}}
\newcommand{\cB}{\mcal{B}}
\newcommand{\cC}{\mcal{C}}
\newcommand{\cD}{\mcal{D}}
\newcommand{\cE}{\mcal{E}}
\newcommand{\cF}{\mcal{F}}
\newcommand{\cG}{\mcal{G}}
\newcommand{\cH}{\mcal{H}}
\newcommand{\cI}{\mcal{I}}
\newcommand{\cJ}{\mcal{J}}
\newcommand{\cK}{\mcal{K}}
\newcommand{\cL}{\mcal{L}}
\newcommand{\cM}{\mcal{M}}
\newcommand{\cN}{\mcal{N}}
\newcommand{\cO}{\mcal{O}}
\newcommand{\cP}{\mcal{P}}
\newcommand{\cQ}{\mcal{Q}}
\newcommand{\cR}{\mcal{R}}
\newcommand{\cS}{\mcal{S}}
\newcommand{\cT}{\mcal{T}}
\newcommand{\cU}{\mcal{U}}
\newcommand{\cV}{\mcal{V}}
\newcommand{\cW}{\mcal{W}}
\newcommand{\cX}{\mcal{X}}
\newcommand{\cY}{\mcal{Y}}
\newcommand{\cZ}{\mcal{Z}}


\newcommand{\scA}{\mathscr{A}}
\newcommand{\scB}{\mathscr{B}}
\newcommand{\scC}{\mathscr{C}}
\newcommand{\scD}{\mathscr{D}}
\newcommand{\scE}{\mathscr{E}}
\newcommand{\scF}{\mathscr{F}}
\newcommand{\scN}{\mathscr{N}}
\newcommand{\scO}{\mathscr{O}}
\newcommand{\scV}{\mathscr{V}}
\newcommand{\scU}{\mathscr{U}}


\newcommand{\ibA}{\mathop{\text{\textit{\textbf{A}}}}}
\newcommand{\ibB}{\mathop{\text{\textit{\textbf{B}}}}}
\newcommand{\ibC}{\mathop{\text{\textit{\textbf{C}}}}}
\newcommand{\ibD}{\mathop{\text{\textit{\textbf{D}}}}}
\newcommand{\ibE}{\mathop{\text{\textit{\textbf{E}}}}}
\newcommand{\ibF}{\mathop{\text{\textit{\textbf{F}}}}}
\newcommand{\ibG}{\mathop{\text{\textit{\textbf{G}}}}}
\newcommand{\ibH}{\mathop{\text{\textit{\textbf{H}}}}}
\newcommand{\ibI}{\mathop{\text{\textit{\textbf{I}}}}}
\newcommand{\ibJ}{\mathop{\text{\textit{\textbf{J}}}}}
\newcommand{\ibK}{\mathop{\text{\textit{\textbf{K}}}}}
\newcommand{\ibL}{\mathop{\text{\textit{\textbf{L}}}}}
\newcommand{\ibM}{\mathop{\text{\textit{\textbf{M}}}}}
\newcommand{\ibN}{\mathop{\text{\textit{\textbf{N}}}}}
\newcommand{\ibO}{\mathop{\text{\textit{\textbf{O}}}}}
\newcommand{\ibP}{\mathop{\text{\textit{\textbf{P}}}}}
\newcommand{\ibQ}{\mathop{\text{\textit{\textbf{Q}}}}}
\newcommand{\ibR}{\mathop{\text{\textit{\textbf{R}}}}}
\newcommand{\ibS}{\mathop{\text{\textit{\textbf{S}}}}}
\newcommand{\ibT}{\mathop{\text{\textit{\textbf{T}}}}}
\newcommand{\ibU}{\mathop{\text{\textit{\textbf{U}}}}}
\newcommand{\ibV}{\mathop{\text{\textit{\textbf{V}}}}}
\newcommand{\ibW}{\mathop{\text{\textit{\textbf{W}}}}}
\newcommand{\ibX}{\mathop{\text{\textit{\textbf{X}}}}}
\newcommand{\ibY}{\mathop{\text{\textit{\textbf{Y}}}}}
\newcommand{\ibZ}{\mathop{\text{\textit{\textbf{Z}}}}}

\newcommand{\ibx}{\mathop{\text{\textit{\textbf{x}}}}}

%\newcommand{\Comp}{\mathop{\mathrm{C}}\nolimits}
%\newcommand{\Komp}{\mathop{\mathrm{K}}\nolimits}
%\newcommand{\Domp}{\mathop{\mathsf{D}}\nolimits}%複体のホモトピー圏
%\newcommand{\Comp}{\mathrm{C}}
%\newcommand{\Komp}{\mathrm{K}}
%\newcommand{\Domp}{\mathsf{D}}%複体のホモトピー圏

\newcommand{\Comp}{\mathsf{C}}
\newcommand{\Komp}{\mathsf{K}}
\newcommand{\Domp}{\mathsf{D}}
\newcommand{\Kompl}{\mathop{\mathsf{K}^\mathrm{+}}\nolimits}
\newcommand{\Kompu}{\mathop{\mathsf{K}^\mathrm{-}}\nolimits}
\newcommand{\Kompb}{\mathop{\mathsf{K}^\mathrm{b}}\nolimits}
\newcommand{\Dompl}{\mathop{\mathsf{D}^\mathrm{+}}\nolimits}
\newcommand{\Dompu}{\mathop{\mathsf{D}^\mathrm{-}}\nolimits}
\newcommand{\Dompb}{\mathop{\mathsf{D}^\mathrm{b}}\nolimits}
\newcommand{\Dompbf}{\mathop{\mathsf{D}_\mathrm{f}^\mathrm{b}}\nolimits}




\newcommand{\CCat}{\Comp(\cat)}
\newcommand{\KCat}{\Komp(\cat)}
\newcommand{\DCat}{\Domp(\cat)}%圏Cの複体のホモトピー圏
\newcommand{\HOM}{\mathop{\mathscr{H}\hspace{-2pt}om}\nolimits}%内部Hom
\newcommand{\RHOM}{\mathop{\mathrm{R}\hspace{-1.5pt}\HOM}\nolimits}

%\newcommand{\muS}{\mathop{\mathrm{SS}}\nolimits}
\newcommand{\muS}[1][]{\mathop{\mathrm{SS}}\nolimits_{#1}}
\newcommand{\RG}{\mathop{\mathrm{R}\hspace{-0pt}\Gamma}\nolimits}
\newcommand{\RHom}{\mathop{\mathrm{R}\hspace{-1.5pt}\Hom}\nolimits}
\newcommand{\Rder}{\mathrm{R}}

\newcommand{\simar}{\mathrel{\overset{\sim}{\rightarrow}}}%同型右矢印
\newcommand{\simarr}{\mathrel{\overset{\sim}{\longrightarrow}}}%同型右矢印
\newcommand{\simra}{\mathrel{\overset{\sim}{\leftarrow}}}%同型左矢印
\newcommand{\simrra}{\mathrel{\overset{\sim}{\longleftarrow}}}%同型左矢印

\newcommand{\hocolim}{{\mathrm{hocolim}}}
\newcommand{\indlim}[1][]{\mathop{\varinjlim}\limits_{#1}}
\newcommand{\sindlim}[1][]{\smash{\mathop{\varinjlim}\limits_{#1}}\,}
\newcommand{\Pro}{\mathrm{Pro}}
\newcommand{\Ind}{\mathrm{Ind}}
\newcommand{\prolim}[1][]{\mathop{\varprojlim}\limits_{#1}}
\newcommand{\sprolim}[1][]{\smash{\mathop{\varprojlim}\limits_{#1}}\,}

\newcommand{\Sh}{\mathrm{Sh}}
\newcommand{\PSh}{\mathrm{PSh}}

\newcommand{\rmD}{\mathrm{D}}

\newcommand{\Lloc}[1][]{\mathord{\mathcal{L}^1_{\mathrm{loc},{#1}}}}
\newcommand{\ori}{\mathord{\mathrm{or}}}
\newcommand{\Db}{\mathord{\mathscr{D}b}}

\newcommand{\codim}{\mathop{\mathrm{codim}}\nolimits}



\newcommand{\gld}{\mathop{\mathrm{gld}}\nolimits}
\newcommand{\wgld}{\mathop{\mathrm{wgld}}\nolimits}


\newcommand{\tens}[1][]{\mathbin{\otimes_{\raise1.5ex\hbox to-.1em{}{#1}}}}
\newcommand{\etens}{\mathbin{\boxtimes}}
\newcommand{\ltens}[1][]{\mathbin{\overset{\mathrm{L}}\tens}_{#1}}
\newcommand{\mtens}[1][]{\mathbin{\overset{\mathrm{\mu}}\tens}_{#1}}
\newcommand{\lltens}[1][]{{\mathop{\tens}\limits^{\mathrm{L}}_{#1}}}
\newcommand{\letens}{\overset{\mathrm{L}}{\etens}}
\newcommand{\detens}{\underline{\etens}}
\newcommand{\ldetens}{\overset{\mathrm{L}}{\underline{\etens}}}
\newcommand{\dtens}[1][]{{\overset{\mathrm{L}}{\underline{\otimes}}}_{#1}}

\newcommand{\blk}{\mathord{\ \cdot\ }}
\newcommand{\mres}[2][]{{\left.{#1}\right\rvert}_{#2}}


%\newcommand{\hocolim}{{\mathrm{hocolim}}}
%\newcommand{\indlim}[1][]{\mathop{\varinjlim}\limits_{#1}}
%\newcommand{\sindlim}[1][]{\smash{\mathop{\varinjlim}\limits_{#1}}\,}
%\newcommand{\Pro}{\mathrm{Pro}}
%\newcommand{\Ind}{\mathrm{Ind}}
%\newcommand{\prolim}[1][]{\mathop{\varprojlim}\limits_{#1}}
%\newcommand{\sprolim}[1][]{\smash{\mathop{\varprojlim}\limits_{#1}}\,}
\newcommand{\proolim}[1][]{\mathop{\text{\rm``{$\varprojlim$}''}}\limits_{#1}}
\newcommand{\sproolim}[1][]{\smash{\mathop{\rm``{\varprojlim}''}\limits_{#1}}}
\newcommand{\inddlim}[1][]{\mathop{\text{\rm``{$\varinjlim$}''}}\limits_{#1}}
\newcommand{\sinddlim}[1][]{\smash{\mathop{\text{\rm``{$\varinjlim$}''}}\limits_{#1}}\,}
\newcommand{\ooplus}{\mathop{\text{\rm``{$\oplus$}''}}\limits}
\newcommand{\bbigsqcup}{\mathop{``\bigsqcup''}\limits}
\newcommand{\bsqcup}{\mathop{``\sqcup''}\limits}
\newcommand{\dsum}[1][]{\mathbin{\oplus_{#1}}}

\newcommand{\Fct}{\mathop{\mathsf{Fct}}\nolimits}
\newcommand{\Red}{\mathop{\mathsf{Red}}\nolimits}
\newcommand{\Cone}{\mathop{\mathsf{Cone}}\nolimits}
\newcommand{\epi}{\mathop{\mathsf{epi}}\nolimits}




\newcommand{\mtan}[1][]{T\hspace{-1.75pt}{#1}}
\newcommand{\mtp}[1][]{{}^{t\!}{#1}} % transpose of the morphism
%\newcommand{\mpb}[1][]{{}^{t\!}{#1}'} % pull-back of the morphism
\newcommand{\mpb}[1][]{{#1}_{\pi}} % pull-back of the morphism
\newcommand{\mdiff}[1][]{{#1}_{d}} % pull-back of the morphism
\newcommand{\mptan}[1][]{{#1}_{\tau}} % pull-back of the morphism


%\newcommand{\pb}[1][]{\mathbin{\mathop{\times}\limits_{#1}}}

\newcommand{\cotb}[1][]{T^{\ast}{#1}}
\newcommand{\conm}[2][]{T_{#1}^{\hspace{0.7pt}\ast}{#2}}
\newcommand{\norb}[2][]{T_{#1}{#2}}

\newcommand{\cotz}[1][]{T^{\ast}_{#1}{#1}}

\newcommand{\cotn}[1][]{\dot{T}^{\ast}{#1}}


\newcommand{\muhom}{\mu hom}
\newcommand{\muHom}[1][]{\mathrm{Hom}^\mu_{\raise1.5ex\hbox to.1em{}#1}}

\newcommand{\mucat}[2][]{\mathsf{D}^{\mathrm{b}}_{#1}(#2)}
%\newcommand{\mucat}[2][]{\mathsf{D}^{\mathrm{b}}_{#1}(#2)}

\newcommand{\conv}[1][]{\mathbin{\mathop{\circ}\limits_{#1}}}

\newcommand{\pb}[1][]{\mathbin{\mathop{\times}\limits_{#1}}}


%================================================
% 自前の定理環境
%   https://mathlandscape.com/latex-amsthm/
% を参考にした
\newtheoremstyle{mystyle}%   % スタイル名
    {5pt}%                   % 上部スペース
    {5pt}%                   % 下部スペース
    {}%              % 本文フォント
    {}%                  % 1行目のインデント量
    {\bfseries}%                      % 見出しフォント
    {.}%                     % 見出し後の句読点
    {12pt}%                     % 見出し後のスペース
    {\thmname{#1}\thmnumber{ #2}\thmnote{{\hspace{2pt}\normalfont (#3)}}}% % 見出しの書式

\theoremstyle{mystyle}
\newtheorem{AXM}{公理}%[section]
\newtheorem{DFN}[AXM]{定義}
\newtheorem{THM}[AXM]{定理}
\newtheorem*{THM*}{定理}
\newtheorem{PRP}[AXM]{命題}
\newtheorem{LMM}[AXM]{補題}
\newtheorem{CRL}[AXM]{系}
\newtheorem{EG}[AXM]{例}
\newtheorem*{EG*}{例}
\newtheorem{RMK}[AXM]{注意}
\newtheorem{CNV}[AXM]{約束}
\newtheorem{CMT}[AXM]{コメント}
\newtheorem*{CMT*}{コメント}
\newtheorem{NTN}[AXM]{記号}
\newtheorem{CLM}[AXM]{Claim}
\newtheorem{Prop}[AXM]{Proposition}

% 定理環境ここまで
%====================================================

\usepackage{framed}
\definecolor{lightgray}{rgb}{0.75,0.75,0.75}
\renewenvironment{leftbar}{%
  \def\FrameCommand{\textcolor{lightgray}{\vrule width 4pt} \hspace{10pt}}% 
  \MakeFramed {\advance\hsize-\width \FrameRestore}}%
{\endMakeFramed}
\newenvironment{redleftbar}{%
  \def\FrameCommand{\textcolor{lightgray}{\vrule width 1pt} \hspace{10pt}}% 
  \MakeFramed {\advance\hsize-\width \FrameRestore}}%
 {\endMakeFramed}



\renewcommand{\Re}{\mathop{\mathrm{Re}}\nolimits}


% =================================





% ---------------------------
% new definition macro
% ---------------------------
% 便利なマクロ集です

% 内積のマクロ
%   例えば \inner<\pphi | \psi> のように用いる
\def\inner<#1>{\langle #1 \rangle}

% ================================
% マクロを追加する場合のスペース 

%=================================







\def\fracinline#1/#2{\mbox{\raise0.5ex\hbox{\footnotesize$#1$}{\hskip-.1em$/$\hskip-.1em}\raise-0.5ex\hbox{\footnotesize$#2$}}}

\begin{document}




\begin{frame}
    \titlepage
\end{frame}
\begin{comment}
    \begin{frame}\frametitle{書くこと}
    \begin{itemize}
        \item 特に興味を持った定理や理論の背景
        \item それを記述するための記号や概念の導入
        \item 正確な主張の紹介
        \item 証明の基本的なアイデアやアウトラインの紹介
        \item または定理の過程を満たす具体例や定理の応用例の紹介
    \end{itemize}
    \end{frame}
\end{comment}
%\setcounter{section}{3}
%\section*{はじめに}
\begin{frame}
    \frametitle{発表内容}
    \tableofcontents
\end{frame}

\section[背景]{背景}
\subsection{層理論}

\begin{frame}
    \frametitle{（前）層の定義}
    \begin{DFN}[層]\label{def:sheaf}
        位相空間\(X\)上の\(\kk\)加群の前層\(F\)とは，
        \(\kk\)加群の族と\(\kk\)加群の射の族の組\[\left(
        \left(F(U)\right)_{U\subset X},
        \left(\rho_{VU}\colon F(U)\to F(V)\right)_{(V\hookrightarrow U)}
        \right)\]で次をみたすもののこと．
        簡単のため\(\kk\)は\(\zz\)か体とする．
        \vspace*{-5pt}
        \begin{enumerate}
            %\renewcommand{\labelenumi}{({\arabic{enumi}})}
            \item $\rho_{UU} = \id_{F(U)}$. 
            \item $W \subset V \subset U$ ならば, 
            $\rho_{WU} = \rho_{WV} \circ \rho_{VU}$.\label{enum:res}
        \end{enumerate}
        \vspace*{-5pt}
        \(F\)が\textbf{層}であるとは，
        \(X\)の任意の開集合\(U\)と
        その開被覆\(\left(U_i\right)_{i\in I}\)に対して，
        次が成り立つこと．
        \vspace*{-5pt}
        \begin{enumerate}[(S1)]
            \setcounter{enumi}{-1}
            \setlength{\leftskip}{10pt}
            \item \(F(\varnothing)=0\).
            \item \(s\in F(U)\)が各\(i\in I\)に対して\(U_i\)上\(s\rvert_{U_{i}}=0\)ならば，
            \(s=0\).\label{sheaf-cond1}
            \item \(\left(s_i\right)_i\in\prod_{i}F(U_i)\)が
            各\(i,j\in I\)で\(U_i\cap U_j\ne\varnothing\)となるものに対して\(U_i\cap U_j\)上\(s_i\rvert_{U_i\cap U_j}-s_j\rvert_{U_i\cap U_j}=0\)ならば，
            \(s\in F(U)\)で各\(U_i\)上\(s\rvert_i=s_i\)となるものが存在する．\label{sheaf-cond2}
        \end{enumerate}
    \end{DFN}
\end{frame}

\begin{frame}
    \frametitle{層の導来圏}
    導来圏の理論を用いると見通しがよい
    \begin{NTN}
        層の有界な複体を対象とし，
        コホモロジー間の同型を誘導する射（擬同型）を可逆にした
        圏を\(\Dompb(\kk_X)\)で表す．
    \end{NTN}
    \(\Dompb(\cc_X)\)では例えば定数層のド・ラーム分解が同型になる：
    \[\vcenter{\xymatrix{
        0\ar[r]&
        \cc_X \ar[r] \ar[d]^-{\mathrm{qis}}&
        0\\
        0\ar[r]&
        \Omega_X^0\ar[r]&
        \Omega_X^1\ar[r]&
        \cdots\ar[r]&
        \Omega_X^n\ar[r]&0
    }}\]
\end{frame}

\begin{frame}
    \frametitle{完全三角}
    導来圏では完全三角 (distinguished triangle) が完全列の代わりになる．
    \[
        F\to G\to H\to+1\quad\mathrm{d.t.}
    \]
    とかいたとき，長完全列
    \begin{align*}
        0&\to H^0(F)\to H^0(G)\to H^0(H)\\
        &\to H^1(F)\to H^1(G)\to H^1(H)\\
        &\to \cdots
    \end{align*}
    が得られる．つまり，上の完全三角で\(H=0\)なら，
    \[
    F\simarr G\quad \mathrm{qis}
    \]となる．
\end{frame}

\subsection{超局所層理論}
\newcommand{\Char}{\mathop{\mathrm{Ch}}\nolimits}
\newcommand{\singsupp}{\mathop{\mathrm{singsupp}}\nolimits}
\newcommand{\SingSpec}{\mathop{\mathrm{S.S.}}\nolimits}
%\newcommand{\SingSpec}{\mathop{\mathrm{WF}_{\mathrm{A}}}\nolimits}

\begin{frame}\frametitle{``超局所''？--- 超局所解析からの復習}
    \[
        \text{``超局所''}=\text{余接束上で局所}
    \]
    いったん，古典的な微分方程式の問題で説明する．
    \begin{itemize}
        \item 余接束：開集合\(\varOmega\subset \rr^n\)に対し
        \(\dot{\pi}\colon\cotn[\varOmega]\coloneqq \cotb[\varOmega]-\cotz[\varOmega]\cong \varOmega\times(\rr^n-\{0\})\to \varOmega\). 
        \item 超関数\(u\)に対し特異台\(
            \singsupp(u)\coloneqq\left\{
                x\in\varOmega;\ \text{\(u(x)\)は\(x\)で解析的でない}
            \right\}\). 
        \item 微分作用素\(P\)に対し\(
            \Char(P)\coloneqq\{(x;\xi)\in\cotn[\varOmega];\ 
                \sigma(P)(x;\xi)=0\}.
            \)\\（ただし\(\sigma(P)\)は\(P\)の``主要表象''）
        \item \(u\)の特異性スペクトル：\(\SingSpec(u)\subset\cotn[\varOmega]\)は\(
            \pi(\SingSpec(u))=\singsupp(u)
        \)をみたす集合．
    \end{itemize}
    それぞれ
    \(\singsupp{u}\subset\varOmega\)
    と
    \(\Char(P), \SingSpec(u)\subset \cotn[\varOmega]\)に注意．
\end{frame}

\begin{frame}\frametitle{超局所解析における定理} 
    \begin{THM}[佐藤の基本定理{\cite{Sato69,Sato70}}]
        \(f\)：既知関数，
        \(P\)：実解析的係数微分作用素とする．
        \(u\)が微分方程式\(Pu=f\)の超関数解のとき
        \[
            \SingSpec(u)\subset \Char(P)\cup\SingSpec(f).
        \]
    \end{THM}
    これを\(f=0\), \(P\)：楕円型の場合に適用すると
    \[
        \Char(P)\cup \SingSpec(f)=\varnothing
    \]
    なので\(\singsupp(u)=\pi(\SingSpec(u))=\varnothing\)．
    すなわち，
    \begin{alertblock}{楕円型正則性定理}
        実解析的係数の楕円型方程式の解は解析的である．
    \end{alertblock}
    このように，底空間の対象を余接束上に持ち上げる（超局所的に考える）と，
    その対象の特異性が分析できる．
\end{frame}

\begin{frame}\frametitle{層の超局所台 --- 定義の動機づけ}
    %\(\Dompb(\kk_X)\ni F\mapsto \SS(F)\subset T^{\ast}X\)
    以上の層理論における類似が超局所層理論である．

    \(F\in\Dompb(\kk_M)\)の超局所台\(\muS(F)\)を定義したい．

    \(p=(x_0;\xi_0)\in\cotb[M]\)とする．
    \(M\)上の実\(C^1\)級関数\(\varphi\)で\(
        d\varphi(x_0)=\xi_0
    \)となるものを考える．

    \(j\)次コホモロジーにおける``方向\(p\)への制限''
    \[
        H^{j}F_{x_{0}}\cong \indlim[U\ni x_{0}]H^{j}(U;F)
        \to
        \indlim[U\ni x_{0}]H^j(U\cap\{x;\varphi(x)<\varphi(x_0)\};F).
    \]
    が同型になるというのが，\(p\)の方向にコホモロジーが伝播するということ．\\
    長完全列を考えると次の定義になる．
\end{frame}


\begin{frame}\frametitle{層の超局所台 --- 定義}
    %\(\Dompb(\kk_X)\ni F\mapsto \SS(F)\subset T^{\ast}X\)    
    \begin{DFN}[層の超局所台]
        \(p=(x_0;\xi_0)\notin \muS(F)\)となるのは，
        \(U\in I_{p}\)で，任意の\(x_1\in X\)と
        \(x_1\)の近傍で定義された，\(
            d\varphi(x_1)\in U
        \)となる\(\varphi\)に対し，\[
            \RG_{\{\varphi\geqq\varphi(x_0)\}}(F)_{x_1}=0
        \]
        となるものが存在するときであるとする．
    \end{DFN}
    超局所台を用いると，関数の臨界点を層係数に一般化することができる．
    \begin{DFN}
        \(\varphi\colon M\to \rr\)の臨界点とは
        \(d\varphi(p)=0\)となる点\(p\in M\)のこと．

        \(F\)に関する\(\varphi\colon M\to \rr\)の臨界点とは
        \(d\varphi(p)\in\muS[](F)\)となる点\(p\in M\)のこと．
    \end{DFN}
\end{frame}

\begin{frame}\frametitle{打切り超局所台 --- 定義}
    \begin{alertblock}{問題点}
        境界値問題など，超局所台では捉えきれない現象がある．
    \end{alertblock}
    
    柏原--Monteiro-Fernandes--Schapira \cite{KFS03a}は，
    このような現象を扱うために
    超局所台を精密化した．
    \begin{DFN}[{\cite[Definition 3.4]{KFS03a}}]
        \(k\in\zz\)とする．
        \(p=(x_0;\xi_0)\notin \muS[k](F)\)となるのは，
        \(p\)の錐状開近傍\(U\)で，任意の\(x_1\in \pi(U)\)と
        \(x_1\)の近傍で定義された，\(
            d\varphi(x_1)\in U
        \)となる\(\varphi\)に対し，任意の次数\(j\leqq k\)で\[
            H^{j}_{\{\varphi\geqq\varphi(x_0)\}}(F)_{x_1}=0
        \]
        となるものが存在するときであるとする．
    \end{DFN}
\end{frame}


\begin{frame}\frametitle{（余）接束の基本図式}
    \(f\colon M\to N\)を多様体の射とする．次の図式が基本的．
    \[\vcenter{\xymatrix@C=22pt@R=26pt{
        TM
        \ar[rd]_-{\tau_M}
        \ar[r]^-{f'}
        &
        M\pb[N] TN
        %\ar@{}[dr]|\square
        \ar[d]_-{}
        \ar[r]^-{\mptan[f]}
        \ar@{}[dr]|\square
        &
        TN
        \ar[d]^-{\tau_N}
        &
        \cotb[M]
        \ar[rd]_-{\pi_M}
        &
        M\pb[N] \cotb[N]
        %\ar@{}[dr]|\square
        \ar[d]_-{}
        \ar[r]^-{\mpb[f]}
        \ar[l]_-{\mdiff[f]}
        \ar@{}[dr]|\square
        &
        \cotb[N]
        \ar[d]^-{\pi_N}
        \\
        &
        M
        \ar[r]^-{f}
        %\ar@{}[dr]|\square
        &
        N,
        &
        &
        M
        \ar[r]^-{f}
        %\ar@{}[dr]|\square
        &
        N.
    }}\]
    \begin{itemize}
        \item \(\mptan[f], \mpb[f]\)は標準射影
        \item \(f'\)は接空間の間の微分写像から定まる射
        \item \(\mdiff[f]\)はその転置から定まる射
    \end{itemize}
\end{frame}
\begin{frame}
    \frametitle{打切り超局所台の関手的性質}
    \begin{PRP}[{\cite[Proposition 5.4.4]{KS90}}]\label{prp:img}
        \(f\colon M\to N\)を多様体の射とし，
        \(F\in\Dompb(\kk_M)\)とする．
        \(f\)が\(\supp(F)\)上固有ならば，
        \[
            \muS[](\Rder{f_{\ast}})
            \subset
            \mpb[f]\mdiff[f]^{-1}\muS[](F)
        \]である．
    \end{PRP}
    打切り版でも同様の事実が成り立つ．
    \begin{PRP}[{\cite[Proposition 4.3]{KFS03a}}]\label{prp:t-img}
        \(f\colon M\to N\)を多様体の射とし，
        \(F\in\Dompb(\kk_M)\)とする．\(k\in\zz\)とする．
        \(f\)が\(\supp(F)\)上固有ならば，
        \[
            \muS[k](\Rder{f_{\ast}})
            \subset
            \mpb[f]\mdiff[f]^{-1}\muS[k](F)
        \]である．
    \end{PRP}
\end{frame}




\section{主結果}
\begin{frame}
    %\frametitle{主結果1}
  \begin{THM}[O., 打切り版非特性変形補題]
    \(X\)をハウスドルフ空間とする．
    \(F\in\Dompl(\kk_X)\)とし，
    \(\left(U_{t}\right)_{t\in\rr}\)を\(X\)の開集合の族とする．
    \(k\in\zz\)とする．
    \(
      Z_{s}\coloneqq \bigcap_{t>s}\overline{U_t- U_s}\ (s\in\rr)
    \)とおく．
    次を仮定する．%\vspace{-8pt}
    \begin{enumerate}[(i)]
      \setlength{\leftskip}{10pt}
      \item \label{item:t-nonchar1}
      任意の\(t\in\rr\)に対し\(U_{t}=\bigcup_{s<t}U_{s}\)である．
      \item \label{item:t-nonchar2}
      任意の実数\(s\leqq t\)に対し\(
        \overline{U_{t}- U_{s}}\cap \supp{F}
      \)はコンパクトである．
      \item \label{item:t-nonchar3}
      任意の実数\(s\in\rr\)と任意の点\(x\in Z_{s}- U_{s}\)に対し，
      \(j\leqq k\)ならば
      \(
        (H^{j}_{X- U_{s}}(F))_{x}=0. 
      \)
      \item \label{item:t-nonchar4} 
      任意の実数\(s< t\)に対し，\(Z_s\subset U_{t}\)が成り立つ．
    \end{enumerate}
    このとき，任意の\(t\in\rr\)に対し，制限射\[
      H^{j}\left(\bigcup_{s\in\rr}U_{s};F\right)
      \to
      H^{j}\left(U_{t};F\right)
    \]は\(j<k\)ならば同型であり，\(j=k\)ならば単射である．
  \end{THM}
\end{frame}

\begin{frame}
    \frametitle{打切り版への拡張}
    \begin{THM}[O., 打切り超局所モースの補題]
        \(k\in\zz\)とする．
        \(\psi\colon M\to \rr\)を\(a\leqq\psi(x)< b\)で
        \(d\psi(x)\notin\muS[k](F)\)となる\(C^1\)級関数で
        \(\supp(F)\)上固有なものとする．
        このとき制限射\[
            H^{j}(M_b;F)\longrightarrow H^{j}(M_a;F)
        \]は\(j\leqq k\)のとき同型であり，\(j=k\)のとき単射である．
    \end{THM}    
\end{frame}


\appendix
\section*{参考文献}

\begin{frame}[allowframebreaks]{参考文献}
  \begin{thebibliography}{90}\beamertemplatetextbibitems
  \bibitem[FM07]{FM07} Teresa Monteiro Fernandes, Ana Rita Martins. 
  \textit{Truncated Microsupport and Hyperbolic Inequalities}, 
  Nagoya Mathematical Journal, 2007, Vol.\ 185, pp.63--91. 
  \bibitem[GKS12]{GKS12} St\'ephane Guillermou, 
  Masaki Kashiwara, Pierre Schapira, 
  \textit{Sheaf Quantization of Hamiltonian Isotopies and Applications to Nondisplaceability Problems}, 
  Duke. Math. J. 161, No. 2, pp.201--245, 2012.  
  \bibitem[KFS03a]{KFS03a} Masaki Kashiwara, Teresa Monteiro Fernandes, Pierre Schapira, 
  \textit{Truncated Microsupport and Holomorphic Solutions 
  of \(D\)-modules}, 
  Ann. Scient. \'Ec. Norm. Sup., s\'erie 4, tome 36, 
  pp.583--399, 2003.
  \bibitem[KFS03b]{KFS03b} Masaki Kashiwara, Teresa Monteiro Fernandes, Pierre Schapira, 
  \textit{Involutivity of Truncated microsupports}, 
  Bull. Soc. math. France, 131 (2), 2003, 
  pp.259--266.
  \bibitem[KS82]{KS82} Masaki Kashiwara and Pierre Schapira,
  \textit{Micro-support des faisceaux: applications 
  aux modules diff{\'e}rentiels,}
  C.~R.~Acad.\ Sci.\ Paris s{\'e}rie I Math 295 8, 487--490 France (1982).
  \bibitem[KS90]{KS90} Masaki Kashiwara, Pierre Schapira, 
  \textit{Sheaves on Manifolds}, 
  Grundlehren der Mathematischen Wissenschaften, 292, Springer, 1990.
  \bibitem[Sato69]{Sato69} Mikio Sato, 
  \textit{Hyperfuncitons and 
  Partial Differential equations}, 
  Proc. Int. Conf. on Functional Analysis, Tokyo, 1969.
  \bibitem[Sato70]{Sato70} Mikio Sato, 
  \textit{Regularity of Hyperfunciton Solutions of 
  Partial Differential equations}, 
  Actes, Congres intern. Math., 
  Tome 2, pp.785--794, 1970.
\end{thebibliography}
\end{frame}


\begin{comment}
\begin{frame}[noframenumbering]
aaa
\end{frame}
\end{comment}

\end{document}